\section{Results I: Candidate Perturbation Screen}
\label{ch:results-screen}

This section reports the completed candidate-versus-native-baseline screen.
It nominates a leading operator but does not yet establish specificity against
generic damage.

\subsection{Nomination and independent circular retest isolated one complete profile}

State persistence $0.95$ was first nominated by an exploratory three-seed pilot
on an earlier circular family without trained distractors. The present screen
retested all candidate strengths on the new distractor-trained circular family
and reused the completed ten-checkpoint N-back grid. Across the 23 non-neutral
shared operator-strength profiles, only persistence $0.95$ again met the
complete descriptive majority pattern (Fig.~\ref{fig:screen}). It showed
slowing with preservation, long-delay selectivity, and distractor selectivity
in four of five circular checkpoints, together with load selectivity in all
ten N-back checkpoints. Relative to the pilot, circular delay selectivity
shrank from $0.290$ to $0.089$ and distractor selectivity from $0.148$ to
$0.015$, while N-back load selectivity remained similar ($0.296$ to $0.250$).
Neutral settings returned zero contrast. Other perturbations often reproduced
one or two components, but not the full combination.

\begin{figure*}[t]
  \centering
  \includegraphics[height=0.70\textheight]{signature_screen.png}
  \caption{Cross-task signature screen.
  Rows are every shared operator-strength setting evaluated in both task
  families; columns show the fraction of independent checkpoint seeds with the
  predicted signature direction. Circular fractions have denominator five and
  N-back fractions denominator ten. State persistence 0.95 was the only row to
  meet the complete majority pattern and additionally reached 4/5 agreement on
  each circular feature with complete N-back agreement. Values quantify
  directional consistency, not participant-level probability or statistical
  significance.}
  \label{fig:screen}
\end{figure*}

\subsection{A 5\% persistence reduction preserved ordinary accuracy while
selectively affecting demanding conditions}

At $p=0.95$, clean delay-20 proportional angular-error change was
$0.049\pm0.092$ across circular checkpoints (95\% $t$-interval
$[-0.065,0.164]$) and remained below $0.20$ in 5/5 seeds. Restricted mean
settling increased by $0.104\pm0.167$ steps ($[-0.103,0.312]$) and was
positive in 4/5 seeds, with a range of $-0.013$ to $0.393$ steps
(Table~\ref{tab:leading}). The circular settling, delay, and distractor
intervals all include zero, so those components are directional majority
tendencies rather than robust across-checkpoint effects. By contrast, N-back
load selectivity was $0.250\pm0.077$ ($[0.195,0.305]$), positive in 10/10
checkpoints, and ranged from $0.110$ to $0.363$. Long-minus-short delay
selectivity was $0.089\pm0.075$ ($[-0.004,0.181]$) and distractor-minus-clean
selectivity was $0.015\pm0.044$ ($[-0.040,0.071]$), each positive in 4/5
checkpoints. Thus the leading perturbation retained the predicted average
ordering, but only the N-back contrast clearly excluded zero under the
seed-level $t$-interval.

\begin{table*}[t]
\caption{State persistence 0.95. Values are mean $\pm$ SD across five
circular or ten N-back checkpoint seeds; intervals are student $t$ 95\%
intervals across those seeds.}
\label{tab:leading}
\centering
\begin{tabular}{lccc}
\toprule
Outcome & Contrast & 95\% interval & Consistency \\
\midrule
Delay-20 error & $0.049\pm0.092$ & $[-0.065,0.164]$ & below 0.20 in 5/5 \\
Settling change (steps) & $0.104\pm0.167$ & $[-0.103,0.312]$ & 4/5 positive \\
Delay selectivity & $0.089\pm0.075$ & $[-0.004,0.181]$ & 4/5 positive \\
Distractor selectivity & $0.015\pm0.044$ & $[-0.040,0.071]$ & 4/5 positive \\
N-back load selectivity & $0.250\pm0.077$ & $[0.195,0.305]$ & 10/10 positive \\
\bottomrule
\end{tabular}
\end{table*}

\subsection{The persistence effect was dose-like but not monotonic across all
features}

Reducing persistence further to 0.90 enlarged mean settling, delay, and N-back
load effects (Fig.~\ref{fig:persistence}). However, slowing with preservation
passed only 1/5 circular checkpoints, so the stronger intervention no longer
expressed the selective accuracy-preserving pattern. Increasing persistence
above one continued to impair N-back preferentially but reversed the circular
delay and distractor contrasts. The joint signature was therefore specific to
a weak reduction around 0.95 rather than monotonic damage at any non-neutral
strength.

\begin{figure*}[t]
  \centering
  \includegraphics[width=0.98\textwidth]{persistence_dose_response.png}
  \caption{Dose-like response to carried-state persistence scaling.
  Grey lines are independently trained checkpoint seeds; blue points and error
  bars are mean $\pm$ SD. The dashed red line marks the leading scale 0.95 and
  the horizontal line marks no change from native performance. A stronger
  reduction increased several impairments but also violated relative
  preservation, whereas increased persistence reversed circular delay and
  distractor selectivity.}
  \label{fig:persistence}
\end{figure*}

\subsection{Alternative operators reproduced partial signatures}

Sensory-input gain 1.20 produced distractor selectivity in 5/5 circular
checkpoints and load selectivity in 9/10 N-back checkpoints, but mean settling
became slightly faster ($-0.010\pm0.023$ steps; positive in only 1/5) and delay
selectivity was positive in only 1/5. Effective time-constant scale 1.10 slowed
settling with preservation in 5/5 ($0.605\pm0.176$ steps) and showed load
selectivity in 10/10, but distractor selectivity was negative in every circular
checkpoint. Conversely, time-constant scale 0.90 showed delay selectivity in
5/5 and distractor selectivity in 4/5 but accelerated settling in every
circular checkpoint ($-0.681\pm0.170$ steps). Synaptic-drive gain 1.05 showed
delay selectivity in 4/5, distractor selectivity in 3/5, and load selectivity
in 10/10, but also accelerated settling in 5/5 ($-0.496\pm0.240$ steps). All of
these near-miss clean delay-20 settling cells were latency-valid; uniqueness
was therefore not an artefact of invalid latency scoring. These near-misses
constrain the result: matching one difficult condition was insufficient to
qualify as the complete pattern, and the discrimination against the conserved
time-constant operator rested on settling direction and distractor selectivity
rather than on large absolute settling magnitude.

The contrast between persistence and the conserved time-constant operator is
mechanistically informative. Both reduced persistence ($p<1$) and reduced
timescale ($s<1$) can weaken retention, but only persistence leaves the drive
coefficient unchanged (Eq.~\ref{eq:persistence}). Reducing the conserved
timescale ($s=0.90$) accelerated post-cue settling while preserving delay and
distractor selectivity; increasing it ($s=1.10$) slowed settling but reversed
distractor selectivity. The joint pattern therefore depended on the asymmetric
carried-state intervention, not merely on any reduction in effective memory
lifetime. Cross-operator comparisons remain uncontrolled for overall clean-task
cost, so dose differences could still contribute to some near-misses.

\subsection{Trained distractor filtering resolved the prior validity failure}

All 945 circular result rows passed the fixation gate, all ten N-back native
checkpoints passed competence, and all $p=0.95$ clean delay-20 cells were
latency-valid. Twelve rows elsewhere in the circular grid were latency-invalid
(all heterogeneous-drive settings with low settled fraction); none of the
near-miss settling comparisons above used those cells. Native distractor mean
error now ranged from $3.46^\circ$ to $4.83^\circ$, and every native distractor
cell had a settled fraction of 1.00. The distractor contrast therefore measured
perturbation of a learned filtering behaviour rather than failure on an
unfamiliar task. As a manipulation check, distractor-window gains 1.10, 1.25,
and 1.50 increased distractor error by 3.7\%, 9.8\%, and 21.5\% on average,
respectively, with positive effects in 5/5 checkpoints. However,
persistence-related distractor selectivity remained small and was negative in
one checkpoint, so the improved validity does not convert the result into a
strong human effect-size match.

A three-seed exploratory pilot had already suggested that persistence $0.95$
can exceed weak Gaussian state noise on N-back load selectivity at nearly
matched 0-back cost (means $0.296$ versus $0.052$ at noise SD $0.025$). That
observation is labelled secondary, used only three N-back checkpoints, and did
not yield an informative circular cost match within the pilot's $0.05$
tolerance; it is therefore not a specificity claim for the present screen.
